\section*{Chapter \ref{ch:ll:the-strategy-engine} --- The Strategy Engine}

\begin{solution}{exo:ll:the-strategy-engine:1}
$34\,200.012345~\text{s}$ is $342\,000\,123$ intervals of \qty{100}{\micro\second}, and $342\,000\,123 \bmod 1\,024 = 507$. The
wheel spans $1\,024 \times \qty{100}{\micro\second} = \qty{102.4}{\milli\second}$, so a timer 5 seconds away is about 48.8
revolutions away: it waits in its slot while the wheel passes it 48 times.
\end{solution}

\begin{solution}{exo:ll:the-strategy-engine:2}
By time, then identifier: (120, 7), (120, 9), (121, 3).
\end{solution}

\begin{solution}{exo:ll:the-strategy-engine:3}
In replay the event's time is when the live engine saw it, so decisions depend on the same numbers. It matches the live run only
if the live engine also took its time from the events (or from a clock whose readings are recorded with the inputs), and if
the recorded inputs are exactly what the live engine received, in the same order.
\end{solution}

\begin{solution}{exo:ll:the-strategy-engine:4}
$k\,(1 - (1 - 1/k)^{100})$ on average. It equals 93 for $k \approx 700$: the outcome depends on several hundred possible
orders (fewer than a hundred would cap the count at $k$; unequal probabilities need more).
\end{solution}

\begin{solution}{exo:ll:the-strategy-engine:5}
The first under the old parameters (two ticks), the second under the new (four): the snapshot is installed before the first
event at or after its effective time. The log shows the version with every action: version 1 before, version 2 after.
\end{solution}

\begin{solution}{exo:ll:the-strategy-engine:6}
Yes: $2 \times 10^6 \times \qty{53}{\nano\second} \approx 0.11$, a utilisation of about 11\%, so no queue builds up (the tail
of \qty{130}{\nano\second} at the p99 does not change that).
\end{solution}

\begin{solution}{exo:ll:the-strategy-engine:7}
Schedule a refresh timer with each quote, due \qty{5}{\second} later, with the quote's identifier in the timer's; when it fires and
the quote is unchanged, emit a replace at the same price and schedule the next refresh. Both engines must use the same
identifiers and the same ordering rule for the replay hashes to agree, and the committed hash is regenerated once from the
C++ engine.
\end{solution}

\begin{solution}{exo:ll:the-strategy-engine:8}
The stamp differs on every run, so the output is never bitwise reproducible, and any code that ever compares it with anything
(a timeout, an age) makes the orders differ too; the system clock can also jump when it is corrected. Take the audit timestamp
from the event's time in the engine, or add it outside the engine where messages leave the machine (chapter~23), and keep it
out of the replayed comparison.
\end{solution}

\begin{solution}{pb:ll:the-strategy-engine:1}
\begin{enumerate}
\item A wall-clock read: 100 distinct outputs; a seeded hash order: 93; two threads on one queue: 49.
\item Every run reads different times, and the times are in the output.
\item The outputs depend on which order the set yields, and the seed maps many seeds to the same order: several hundred orders,
repeated now and then in a hundred runs.
\item Many interleavings end in the same output when the threads happen to alternate as they did before; more contention or
slower consumers would give more.
\item Log the output of both runs with sequence numbers and compare them record by record (or bisect on the hash of prefixes).
\item The wall clock: identical decisions with different time stamps; the hash order: the same actions in a different order or
with different identifiers; the threads: a book state that the input never produced.
\item The seed is drawn at start-up, once per process: reruns inside one process share it.
\item The thread interleaving, which a debugger changes by stopping threads.
\item Take time from the events through an \omterm{def:ll:the-strategy-engine:wheel}{injected clock}, and stamp outside the engine.
\item Keep hash maps for lookups only, and iterate in a defined order (an array of the keys, an ordered map, or a sort) when the
order decides anything.
\item Merge the two inputs by a rule (time, then source, then sequence) in one thread before the engine, or give each input its
own ring and let the engine merge.
\item One; the firm's engine gave one.
\item \textbf{Named result.} In a hundred reruns of the same 6\,000 events: 100 distinct outcomes with a wall-clock read, 93 with
a seeded hash order, 49 with two threads on one queue, 100 with all three, and one after the fixes (and one for the firm's
engine).
\item A backtest that is not reproducible cannot be debugged or compared: a small difference can hide a large bug, and two
versions of a strategy cannot be told apart from noise.
\item Floating-point reductions in a different order, different compiler flags or libraries (chapter~8), different input data
(a missed packet in a recording), and uninitialised memory.
\item Integer operations are exact and defined the same way in both languages, while floating-point results depend on the order
and contraction of operations.
\item Replay a recorded day twice and compare hashes; replay it on each build of the engine and compare with the committed
hash; run the replay under a sanitiser.
\item Any incident can be replayed exactly on a desk, with a debugger, as often as needed.
\item Outside the decision path: monitoring, logging of timings, operational tools whose outputs are not part of the replayed
state.
\item One thread, ordered inputs, time taken from the events, and nothing iterated in an order that varies.
\end{enumerate}
\end{solution}

\begin{solution}{iq:ll:the-strategy-engine:1}
No locks, no coordination, the state in one core's caches, and a single order of events that makes the engine deterministic
and replayable. The cost: one core's throughput is the limit, a slow handler delays every other event, and work that could run
in parallel (analytics, logging) must be moved to other threads through rings.
\iqlookfor{determinism and locality against a single core's capacity.}
\end{solution}

\begin{solution}{iq:ll:the-strategy-engine:2}
A hashed timing wheel: constant-time scheduling and cancellation, slots visited as time advances, a hierarchical wheel or
revolution counts for long timers; fire due timers in a defined order (time, then identifier).
\iqlookfor{the wheel, and ordering for determinism.}
\end{solution}

\begin{solution}{iq:ll:the-strategy-engine:3}
Inputs first (are the recorded events what production received, in the same order?), then time (does any code read the clock
instead of the event's time?), then order (threads, hash iteration), then numerics (floating-point order, flags), then
parameters (the versions in force). Compare the two outputs record by record to find the first difference.
\iqlookfor{inputs, time, order, numerics, parameters.}
\end{solution}

\begin{solution}{iq:ll:the-strategy-engine:4}
Publish complete, versioned snapshots on a control ring with an effective time; the engine installs one between events, never
during one, and stamps every action with the version in force; in replay the snapshots are inputs.
\iqlookfor{whole snapshots between events, and versioning.}
\end{solution}

\begin{solution}{iq:ll:the-strategy-engine:5}
Allocation, system calls, locks, waiting, text formatting, unbounded loops and anything that varies from run to run. Check with
allocation counters in tests, tracing of system calls, replay hashes, and measured tail latencies per event.
\iqlookfor{a list and a way to check each item.}
\end{solution}

\begin{solution}{iq:ll:the-strategy-engine:6}
One thread per engine fed by rings; inputs merged by a rule and recorded; time injected (simulated in replay, real live); timers
on a wheel with ordered firing; parameters as versioned snapshots between events; no hash-order or floating-point dependence
in decisions; enforced by replay tests with committed hashes, in continuous integration and after every release.
\iqlookfor{rules plus tests that enforce them.}
\end{solution}
